\SourceLeaf{262}
\[
\begin{aligned}
(1,2)&=aa'\mathfrak{n}''\\
(1,3)&=aa''\mathfrak{n}'\\
(1,4)&=ab'\mathfrak{n}''+ab''\mathfrak{n}'\\
(1,5)&=a'a''\mathfrak{n}\\
(1,6)&=a'b\mathfrak{n}''+a'b''\mathfrak{n}\\
(1,7)&=a''b\mathfrak{n}'+a''b'\mathfrak{n}\\
(1,8)&=bb'\mathfrak{n}''+bb''\mathfrak{n}'+b'b''\mathfrak{n}+\mathfrak{D}\mathfrak{n}\mathfrak{n}'\mathfrak{n}'',\\
(2,3)&=ab''\mathfrak{n}'-ab'\mathfrak{n}''\\
(2,4)&=ac''\mathfrak{n}'\\
(2,5)&=a'b''\mathfrak{n}-a'b\mathfrak{n}''\\
(2,6)&=a'c''\mathfrak{n}\\
(2,7)&=bb''\mathfrak{n}'+b'b''\mathfrak{n}-bb'\mathfrak{n}''-\mathfrak{D}\mathfrak{n}\mathfrak{n}'\mathfrak{n}'',\\
(2,8)&=bc''\mathfrak{n}'+b'c''\mathfrak{n}\\
(3,4)&=ac'\mathfrak{n}''
\end{aligned}
\]
\[
\begin{aligned}
(3,5)&=a''b'\mathfrak{n}-a''b\mathfrak{n}'\\
(3,6)&=bb'\mathfrak{n}''+b'b''\mathfrak{n}-bb''\mathfrak{n}'-\mathfrak{D}\mathfrak{n}\mathfrak{n}'\mathfrak{n}''\\
(3,7)&=a''c'\mathfrak{n}\\
(3,8)&=bc'\mathfrak{n}''+b''c\mathfrak{n}\\
(4,5)&=b'b''\mathfrak{n}-bb'\mathfrak{n}''-bb''\mathfrak{n}'+\mathfrak{D}\mathfrak{n}\mathfrak{n}'\mathfrak{n}''\\
(4,6)&=b'c''\mathfrak{n}-bc''\mathfrak{n}'\\
(4,7)&=b''c'\mathfrak{n}-bc'\mathfrak{n}''\\
(4,8)&=c'c''\mathfrak{n}\\
(5,6)&=ca'\mathfrak{n}''\\
(5,7)&=ca''\mathfrak{n}'\\
(5,8)&=b'c\mathfrak{n}''+b''c\mathfrak{n}'\\
(6,7)&=b''c\mathfrak{n}'-b'c\mathfrak{n}''\\
(6,8)&=cc''\mathfrak{n}'\\
(7,8)&=cc'\mathfrak{n}''
\end{aligned}
\]
quas per $\Phi$ designabimus, novemque aliae:
\[
\begin{aligned}
(10)(11)-(9)(12)&=a\mathfrak{n}'\mathfrak{n}''\mathfrak{A}\\
(1)(12)-(2)(11)-(3)(10)+(4)(9)&=2a\mathfrak{n}'\mathfrak{n}''\mathfrak{B}\\
(2)(3)-(1)(4)&=a\mathfrak{n}'\mathfrak{n}''\mathfrak{C}\\
-(9)(16)+(10)(15)+(11)(14)-(12)(13)&=2b\mathfrak{n}'\mathfrak{n}''\mathfrak{A}\\
\begin{gathered}
(1)(16)-(2)(15)-(3)(14)+(4)(13)\{\\
+(5)(12)-(6)(11)-(7)(10)+(8)(9)\}
\end{gathered}&=4b\mathfrak{n}'\mathfrak{n}''\mathfrak{B}\\
-(1)(8)+(2)(7)+(3)(6)-(4)(5)&=2b\mathfrak{n}'\mathfrak{n}''\mathfrak{C}\\
(14)(15)-(13)(16)&=c\mathfrak{n}'\mathfrak{n}''\mathfrak{A}\\
(5)(16)-(6)(15)-(7)(14)+(8)(13)&=2c\mathfrak{n}'\mathfrak{n}''\mathfrak{B}\\
(6)(7)-(5)(8)&=c\mathfrak{n}'\mathfrak{n}''\mathfrak{C}
\end{aligned}
\]
quas designabimus per $\Psi$\HistoricalNote{*)}{Observare convenit, 18 alias aequationes his $\Psi$ similes erui posse, in quibus ad dextram loco factorum $a$, $2b$, $c$ habeantur $a'$, $2b'$, $c'$; $a''$, $2b''$, $c''$: sed has quum ad institutum nostrum non sint necessariae, omittimus.}.

IV. Originem omnium harum 37 aequationum deducere nimis prolixum foret: sufficiet quasdam confirmavisse, ad quarum instar reliquae haud difficulter demonstrari poterunt.

\leavevmode\SourceLeaf{263}
1) Habetur
\[
\begin{aligned}
(1,2)&=(1)(10)-(2)(9)\\
&=(\mathfrak{p}\mathfrak{q}'-\mathfrak{q}\mathfrak{p}')pp\\&\quad+(\mathfrak{p}\mathfrak{q}'''-\mathfrak{q}\mathfrak{p}'''-\mathfrak{p}'\mathfrak{q}''+\mathfrak{q}'\mathfrak{p}'')pq\\&\quad+(\mathfrak{p}''\mathfrak{q}'''-\mathfrak{q}''\mathfrak{p}''')qq\\
&=\mathfrak{n}''(App+2Bpq+Cqq)=\mathfrak{n}''aa'
\end{aligned}
\]
quae est aequ. prima.

2) Fit
\[
\begin{aligned}(1,3)&=(1)(11)-(3)(9)\\&=(\mathfrak{p}\mathfrak{q}''-\mathfrak{q}\mathfrak{p}'')(pq'-qp')\\&=a''\mathfrak{N}an'=aa''\mathfrak{n}'\end{aligned}
\]
aequ. secunda.

3) Erit
\[
\begin{aligned}
(1,8)&=(1)(16)-(8)(9)\\
&=(\mathfrak{p}\mathfrak{q}'-\mathfrak{q}\mathfrak{p}')pp'''\\&\quad+(\mathfrak{p}\mathfrak{q}'''-\mathfrak{q}\mathfrak{p}''')pq'''\\&\quad-(\mathfrak{p}'\mathfrak{q}''-\mathfrak{q}'\mathfrak{p}'')qp'''\\&\quad+(\mathfrak{p}''\mathfrak{q}'''-\mathfrak{q}''\mathfrak{p}''')qq'''\\
&=\mathfrak{n}''(App'''+B(pq'''+qp''')+Cqq''')\\&\quad+b''\mathfrak{N}(pq'''-qp''')\\
&=\mathfrak{n}''(bb'+\surd dd')+b''\mathfrak{N}(b'n+bn')\text{\textsuperscript{*)}}\\
&=\mathfrak{n}''bb'+\mathfrak{n}'bb''+\mathfrak{n}b'b''+\mathfrak{D}\mathfrak{n}\mathfrak{n}'\mathfrak{n}''
\end{aligned}
\]
\HistoricalDisplayNote{*)}{Hoc sequitur ex aequ. 10 art. 235 et sqq. Quantitas radicalis $\surd dd'$ fit $=Dnn'=\mathfrak{D}nn'\mathfrak{N}\mathfrak{N}=\mathfrak{D}\mathfrak{n}\mathfrak{n}'$.}
aequatio octava in $\Phi$. Aequationes reliquas lectoribus confirmandas linquimus.

V. Ex aequatt. $\Phi$ sequitur, viginti octo numeros $(1,2)$, $(1,3)$ etc. nullum divisorem communem habere, sequenti modo. Primo observamus, viginti septem producta e ternis factoribus. quorum primus vel $\mathfrak{n}$, secundus aliquis numerorum $a'$, $2b'$, $c'$, tertiusque aliquis numerorum $a''$, $2b''$, $c''$; vel primus $\mathfrak{n}'$, secundus aliquis e numeris $a$, $2b$, $c$, tertius aliquis numerorum $a''$, $2b''$, $c''$; vel denique primus $\mathfrak{n}''$, secundus aliquis numerorum $a$, $2b$, $c$ tertiusque aliquis e numeris $a'$, $2b'$, $c'$ --- singula haec viginti septem producta propter aequatt. $\Phi$ aequalia esse vel alicui ex viginti octo numeris $(1,2)$, $(1,3)$ etc. vel plurium summae aut differentiae (\textit{e.g.}\ $\mathfrak{n}a'a''=(1,5)$, $2\mathfrak{n}a'b''=(1,6)+(2,5)$, $4\mathfrak{n}b'b''=(1,8)+(2,7)+(3,6)+(4,5)$, et sic de reliquis); quamobrem si hi numeri divisorem communem haberent, hic necessario etiam omnia illa producta metiri deberet. Hinc vero facile deducitur adiumento art. 40 et per methodum saepius in praecedentibus adhibitam, eundem divisorem etiam numeros $\mathfrak{n}m'm''$, $\mathfrak{n}'mm''$, $\mathfrak{n}''mm'$ metiri debere, adeoque horum
\SourceLeaf{264}
quadrata quae sunt $\frac{dm'm'm''m''}{\mathfrak{D}}$, $\frac{d'mmm''m''}{\mathfrak{D}}$, $\frac{d''mmm'm'}{\mathfrak{D}}$ per illius quadratum divisibilia esse, \textit{Q. E. A.}, quoniam per I trium numeratorum divisor communis maximus est $\mathfrak{D}$, adeoque quadrata ipsa divisorem communem habere nequeunt.

VI. Haec omnia pertinent ad transformationem formae $\mathfrak{F}$ in $ff'f''$; et ex transformationibus formae $F$ in $ff'$ formaeque $\mathfrak{F}$ in $Ff''$ deducta sunt. Sed prorsus simili modo e transformationibus formae $F'$ in $ff''$ formaeque $\mathfrak{F}'$ in $F'f'$ derivabitur transformatio formae $\mathfrak{F}'$ in $ff'f''$ talis:
\[
\begin{aligned}
\mathfrak{X}'&=(1)'xx'x''+(2)'xx'y''+(3)'xy'x''+\text{etc.}\\
\mathfrak{Y}'&=(9)'xx'x''+(10)'xx'y''+(11)'xy'x''+\text{etc.}
\end{aligned}
\]
(designando omnes coefficientes similiter ut in transformatione formae $\mathfrak{F}$ in $ff'f''$, singulisque distinctionis caussa lineolam affigendo), ex qua perinde ut ante 28 aequationes ipsis $\Phi$ analogae deducentur, quas per $\Phi'$ designabimus, novemque aliae ipsis $\Psi$ analogae, quas exprimemus per $\Psi'$. Scilicet denotando
\[
(1)'(10)'-(2)'(9)'\text{ per }(1,2)',\quad
(1)'(11)'-(3)'(9)'\text{ per }(1,3)'
\]
aequationes $\Phi'$ erunt
\[
(1,2)'=aa'\mathfrak{n}'',\quad (1,3)'=aa''\mathfrak{n}',\quad\text{etc.}
\]
aequationes $\Psi'$ autem
\[
(10)'(11)'-(9)'(12)'=a\mathfrak{n}'\mathfrak{n}''\mathfrak{A}',\quad\text{etc.}
\]
(Evolutionem uberiorem brevitatis gratia lectoribus relinquimus; ceterum periti novum calculum ne necessarium quidem esse, sed analysin primam per analogiam facile huc transferri posse invenient). Quibus ita factis, ex $\Phi$ et $\Phi'$ statim sequitur
\[
(1,2)=(1,2)',\quad (1,3)=(1,3)',\quad (1,4)=(1,4)',\quad (2,3)=(2,3)',\quad\text{etc.}
\]
hinc vero et inde quod omnes $(1,2)$, $(1,3)$, $(2,3)$ etc. divisorem communem (per V) non habent, adiumento lemmatis art. 234 concluditur, quatuor numeros integros $\alpha$, $\beta$, $\gamma$, $\delta$ ita determinari posse, ut fiat
\[
\begin{aligned}
\alpha(1)'+\beta(9)'&=(1),\\\alpha(2)'+\beta(10)'&=(2),\\\alpha(3)'+\beta(11)'&=(3),\quad\text{etc.}\\
\gamma(1)'+\delta(9)'&=(9),\\\gamma(2)'+\delta(10)'&=(10),\\\gamma(3)'+\delta(11)'&=(11),\quad\text{etc.}
\end{aligned}
\]
atque $\alpha\delta-\beta\gamma=1$.

\SourceLeaf{265}
VII. Hinc atque substituendo ex tribus aequatt. primis $\Psi$ valores ipsorum $a\mathfrak{A}$, $a\mathfrak{B}$, $a\mathfrak{C}$, et ex tribus aequ. primis $\Psi'$ valores ipsorum $a\mathfrak{A}'$, $a\mathfrak{B}'$, $a\mathfrak{C}'$ facile confirmatur fore
\[
\begin{aligned}
a(\mathfrak{A}\alpha\alpha+2\mathfrak{B}\alpha\gamma+\mathfrak{C}\gamma\gamma)&=a\mathfrak{A}'\\
a(\mathfrak{A}\alpha\beta+\mathfrak{B}(\alpha\delta+\beta\gamma)+\mathfrak{C}\gamma\delta)&=a\mathfrak{B}'\\
a(\mathfrak{A}\beta\beta+2\mathfrak{B}\beta\delta+\mathfrak{C}\delta\delta)&=a\mathfrak{C}'
\end{aligned}
\]
unde, nisi $a=0$, manifesto sequitur, formam $\mathfrak{F}$ transire in $\mathfrak{F}'$ per substitutionem propriam $\alpha$, $\beta$, $\gamma$, $\delta$. --- Adhibendo autem loco trium aequationum primarum in $\Psi$ et $\Psi'$ tres sequentes, facile confirmabuntur tres aequationes modo traditis omnino similes, in quibus loco factoris $a$ ubique invenitur $b$; unde patet, eandem conclusionem etiamnum valere, si modo non sit $b=0$. Denique adhibendo tres ultimas aequationes $\Psi$, $\Psi'$ invenietur eodem modo, conclusionem veram esse, nisi $c=0$. Quocirca, quum certo omnes $a$, $b$, $c$ simul $=0$ esse nequeant, necessario forma $\mathfrak{F}$ per subst. $\alpha$, $\beta$, $\gamma$, $\delta$ transibit in $\mathfrak{F}'$, adeoque huic formae proprie aequivalebit. \textit{Q. E. D.}

\Needspace{8\baselineskip}
\Article{241}
Talem formam ut $\mathfrak{F}$ vel $\mathfrak{F}'$, quae oritur, si una trium formarum datarum componitur cum ea quae ex compositione duarum reliquarum resultat, \textit{ex his tribus formis compositam} vocabimus. patetque ex art. praec., nihil hic interesse, quonam ordine tres formae componantur. Simili modo propositis quotcunque formis $f$, $f'$, $f''$, $f'''$ etc. (quarum determinantes rationem quadratorum inter se habere debent), si forma $f$ componitur cum $f'$, resultans cum $f''$, quae hinc oritur cum $f'''$ etc.: forma quae ad finem huius operationis prodit \textit{ex omnibus formis} $f$, $f'$, $f''$, $f'''$ etc. composita dicetur. Facile vero demonstratur, etiam hic arbitrarium esse, quonam ordine formae componantur; \textit{i. e.}\ quocunque ordine hae formae componantur, formas ex compositione oriundas semper proprie aequivalentes esse. --- Porro manifestum est, si formis $f$, $f'$, $f''$ etc. proprie aequivaleant formae $g$, $g'$, $g''$ etc. resp., formam compositam ex his proprie aequivalentem fore formae ex illis compositae.

\SourceLeaf{266}
\Needspace{8\baselineskip}
\Article{242}
Propositiones praecedentes formarum compositionem maxima universalitate complectuntur; progredimur iam ad applicationes magis particulares, per quas illarum ordinem interrumpere noluimus. Ac primo quidem resumemus problema art. 236, quod per conditiones sequentes limitabimus: \textit{primo} ut formae componendae eundem determinantem habeant, sive sit $d=d'$; \textit{secundo} ut $m$. $m'$ sint inter se primi; \textit{tertio} ut forma quaesita directe ex utraque $f$, $f'$ composita sit. Hinc etiam $mm$, $m'm'$ inter se primi erunt; quare divisor communis maximus numerorum $dm'm'$, $d'mm$ \textit{i. e.}\ $D$ fiet $=d=d'$, atque $n=n'=1$. Quatuor quantitates $\mathfrak{Q}$, $\mathfrak{Q}'$, $\mathfrak{Q}''$, $\mathfrak{Q}'''$, quae ad libitum assumi possunt, statuemus $=-1$, $0$, $0$, $0$ resp., quod semper licebit unico casu excepto, ubi $a$, $a'$, $b+b'$ simul sunt $=0$, ad quem igitur hic non respiciemus; manifesto autem hic casus occurrere nequit nisi in formis determinantis positivi quadrati. Tunc patet, $\mu$ fieri divisorem communem maximum numerorum $a$, $a'$, $b+b'$; numeros $\mathfrak{P}'$, $\mathfrak{P}''$, $\mathfrak{P}'''$ ita accipi debere ut fiat
\[
\mathfrak{P}'a+\mathfrak{P}''a'+\mathfrak{P}'''(b+b')=\mu
\]
ipsum $\mathfrak{P}$ vero omnino arbitrarium esse. Hinc provenit, substituendo l. c. pro $p$, $q$, $p'$, $q'$ etc. valores suos:
\[
\begin{gathered}A=\frac{aa'}{\mu\mu},\\ B=\frac{1}{\mu}(\mathfrak{P}aa'+\mathfrak{P}'ab'+\mathfrak{P}''a'b+\mathfrak{P}'''(bb'+D))\end{gathered}
\]
$C$ autem per aequationem $AC=BB-D$ poterit determinari, si modo non simul $a$ et $a'=0$.

In hac igitur solutione valor ipsius $A$ non pendet a valoribus ipsorum $\mathfrak{P}$, $\mathfrak{P}'$, $\mathfrak{P}''$, $\mathfrak{P}'''$ (qui infinitis modis diversis determinari possunt); $B$ autem alios valores obtinebit tribuendo his numeris alios valores, operaeque pretium est investigare, quomodo omnes valores ipsius $B$ inter se connexi sint. Ad hunc finem observamus

I. Quomodocunque determinentur $\mathfrak{P}$, $\mathfrak{P}'$, $\mathfrak{P}''$, $\mathfrak{P}'''$, valores ipsius $B$ inde prodeuntes omnes congruos esse secundum modulum $A$. Ponamus, si statuatur
\[
\begin{gathered}\mathfrak{P}=\mathfrak{p},\\ \mathfrak{P}'=\mathfrak{p}',\\ \mathfrak{P}''=\mathfrak{p}'',\\ \mathfrak{P}'''=\mathfrak{p}''',\\\text{fieri } B=\mathfrak{B}\end{gathered}
\]
faciendo autem
\[
\begin{gathered}\mathfrak{P}=\mathfrak{p}+\mathfrak{d},\\ \mathfrak{P}'=\mathfrak{p}'+\mathfrak{d}',\\ \mathfrak{P}''=\mathfrak{p}''+\mathfrak{d}''.\\ \mathfrak{P}'''=\mathfrak{p}'''+\mathfrak{d}''',\\\text{prodire } B=\mathfrak{B}+\mathfrak{D}\end{gathered}
\]
\SourceLeaf{267}
Tunc igitur erit
\[
\begin{gathered}a\mathfrak{d}'+a'\mathfrak{d}''+(b+b')\mathfrak{d}'''=0,\\
aa'\mathfrak{d}+ab'\mathfrak{d}'+a'b\mathfrak{d}''+(bb'+D)\mathfrak{d}'''=\mu\mathfrak{D}\end{gathered}
\]
Multiplicando aequationis posterioris partem primam per $a\mathfrak{p}'+a'\mathfrak{p}''+(b+b')\mathfrak{p}'''$, secundam per $\mu$, et subtrahendo a producto primo quantitatem
\[
(ab'\mathfrak{p}'+a'b\mathfrak{p}''+(bb'+D)\mathfrak{p}''')(a\mathfrak{d}'+a'\mathfrak{d}''+(b+b')\mathfrak{d}''')
\]
quae propter aequationem priorem manifesto erit $=0$, habebitur evolutione facta et sublatis quae se destruunt
\[
\begin{aligned}aa'(&\mu\mathfrak{d}\\&+((b'-b)\mathfrak{p}''+c'\mathfrak{p}''')\mathfrak{d}'\\&+((b-b')\mathfrak{p}'+c\mathfrak{p}''')\mathfrak{d}''\\&-(c'\mathfrak{p}'+c\mathfrak{p}'')\mathfrak{d}''')=\mu\mu\mathfrak{D}\end{aligned}
\]
unde manifesto $\mu\mu\mathfrak{D}$ per $aa'$, sive $\mathfrak{D}$ per $\frac{aa'}{\mu\mu}$ \textit{i. e.}\ per $A$ divisibilis erit, atque
\[
\mathfrak{B}\equiv\mathfrak{B}+\mathfrak{D}\ (\mathrm{mod.}\ A)
\]

II. Si valores $\mathfrak{p}$, $\mathfrak{p}'$, $\mathfrak{p}''$ $\mathfrak{p}'''$ ipsorum $\mathfrak{P}$, $\mathfrak{P}'$, $\mathfrak{P}''$, $\mathfrak{P}'''$ reddant $B=\mathfrak{B}$, inveniri posse alios valores horum numerorum, ex quibus $B$ nanciscatur valorem quemcunque datum ipsi $\mathfrak{B}$ secundum mod. $A$ congruum, puta $\mathfrak{B}+kA$. Primo observamus, quatuor numeros $\mu$, $c$, $c'$, $b-b'$ divisorem communem habere non posse; nam si quem haberent, hic metiretur sex numeros $a$, $a'$, $b+b'$, $c$, $c'$, $b-b'$ adeoque tum ipsos $a$, $2b$, $c$, tum ipsos $a'$, $2b'$, $c'$ et proin etiam ipsos $m$, $m'$, qui per hyp. inter se sunt primi. Quamobrem quatuor numeri integri $h$, $h'$, $h''$, $h'''$ poterunt assignari tales ut fiat
\[
h\mu+h'c+h''c'+h'''(b-b')=1
\]
Quo facto si statuitur
\[
\begin{aligned}
kh&=\mathfrak{d},\\k(h''(b+b')-h'''a')&=\mu\mathfrak{d}'\\
k(h'(b+b')+h'''a)&=\mu\mathfrak{d}'',\\-k(h'a'+h''a)&=\mu\mathfrak{d}'''
\end{aligned}
\]
patet, ipsos $\mathfrak{d}$, $\mathfrak{d}'$, $\mathfrak{d}''$, $\mathfrak{d}'''$ esse integros; porro facile confirmatur, fieri
\[
a\mathfrak{d}'+a'\mathfrak{d}''+(b+b')\mathfrak{d}'''=0
\]
\[
\begin{aligned}aa'\mathfrak{d}+ab'\mathfrak{d}'+a'b\mathfrak{d}''+(bb'+D)\mathfrak{d}'''\\&=\frac{aa'k}{\mu}(\mu h+ch'+c'h''+(b-b')h''')\\&=\mu kA\end{aligned}
\]
Ex aequatione priori patet, etiam $\mathfrak{p}+\mathfrak{d}$, $\mathfrak{p}'+\mathfrak{d}'$, $\mathfrak{p}''+\mathfrak{d}''$, $\mathfrak{p}'''+\mathfrak{d}'''$ esse valores ipsorum $\mathfrak{P}$, $\mathfrak{P}'$, $\mathfrak{P}''$, $\mathfrak{P}'''$; ex posteriori, hos valores producere $B=\mathfrak{B}+kA$. \textit{Q. E. D.} ---
\SourceLeaf{268}
Hinc perspicuum est, $B$ semper ita determinari posse ut iaceat inter $0$ et $A-1$ incl., siquidem $A$ est positivus; vel inter $0$ et $-A-1$ si $A$ negativus.

\Needspace{8\baselineskip}
\Article{243}
Ex aequationibus
\[
\begin{gathered}\mathfrak{P}'a+\mathfrak{P}''a'+\mathfrak{P}'''(b+b')=\mu,\\
B=\frac{1}{\mu}(\mathfrak{P}aa'+\mathfrak{P}'ab'+\mathfrak{P}''a'b+\mathfrak{P}'''(bb'+D))\end{gathered}
\]
deducitur
\[
\begin{aligned}B&=b+\frac{a}{\mu}(\mathfrak{P}a'+\mathfrak{P}'(b'-b)-\mathfrak{P}'''c)\\&=b'+\frac{a'}{\mu}(\mathfrak{P}a+\mathfrak{P}''(b-b')-\mathfrak{P}'''c')\end{aligned}
\]
quare
\[
B\equiv b\ (\mathrm{mod.}\ \tfrac{a}{\mu})\quad\text{et}\quad B\equiv b'\ (\mathrm{mod.}\ \tfrac{a'}{\mu})
\]
Quoties $\frac{a}{\mu}$, $\frac{a'}{\mu}$ inter se primi sunt, inter $0$ et $A-1$ (sive inter $0$ et $-A-1$ quando $A$ est negativus) unicus tantum numerus iacebit qui secundum mod. $\frac{a}{\mu}$ sit $\equiv b$, et $\equiv b'$ sec. mod. $\frac{a'}{\mu}$; qui si statuitur $=B$ atque $\frac{BB-D}{A}=C$, palam est, $(A,B,C)$ e formis $(a,b,c)$, $(a',b',c')$ compositam fore. In hoc itaque casu ad inventionem formae compositae ad numeros $\mathfrak{P}$, $\mathfrak{P}'$, $\mathfrak{P}''$, $\mathfrak{P}'''$ non amplius oportet respicere\HistoricalNote{*)}{Quod semper efficitur adhibendo congruentias
\[
\frac{aB}{\mu}\equiv\frac{ab'}{\mu},\quad\frac{a'B}{\mu}\equiv\frac{a'b}{\mu},\quad\frac{(b+b')B}{\mu}\equiv\frac{bb'+D}{\mu}\ (\mathrm{mod.}\ A)
\]
}. Ita \textit{e. g.}\ si quaeritur forma e formis $(10,3,11)$, $(15,2,7)$ composita, erunt $a$, $a'$, $b+b'$ resp. $=10,15,5$; $\mu=5$; hinc $A=6$; $B\equiv3\ (\mathrm{mod.}\ 2)$ et $\equiv2\ (\mathrm{mod.}\ 3)$, unde $B=5$ atque $(6,5,21)$ forma quaesita. --- Ceterum conditio ut $\frac{a}{\mu}$, $\frac{a'}{\mu}$ inter se primi sint, omnino aequivalet huic, ut numeri duo $a$, $a'$ divisorem communem maiorem non habeant quam tres $a$, $a'$, $b+b'$, sive, quod eodem redit, ut divisor communis maximus numerorum $a$, $a'$ etiam numerum $b+b'$ metiatur. Notentur imprimis sequentes casus particulares:

1) Propositis duabus formis $(a,b,c)$, $(a',b',c')$ eiusdem determinantis $D$ ita comparatis ut divisor comm. max. numerorum $a$, $2b$, $c$ primus sit ad div. comm. max. num. $a'$, $2b'$, $c'$, atque $a$ primus ad $a'$: forma ex his composita $(A,B,C)$ invenitur faciendo $A=aa'$, $B\equiv b\ (\mathrm{mod.}\ a)$ et $\equiv b'\ (\mathrm{mod.}\ a')$, $C=\frac{BB-D}{A}$. Hic casus semper locum habet, quando altera formarum componendarum est forma principalis, puta $a=1$, $b=0$, $c=-D$. Tunc erit $A=a'$, $B$ statui poterit
\SourceLeaf{269}
$=b'$, unde fiet $C=c'$; quare \textit{ex forma principali et quacunque alia forma eiusdem determinantis composita est haec forma ipsa.}

2) Si duae formae \textit{oppositae} proprie primitivae sunt componendae, puta $(a,b,c)$ et $(a,-b,c)$, erit $\mu=a$. Hinc facile perspicitur, formam principalem $(1,0,-D)$ ex illis esse compositam.

3) Propositis quotcunque formis proprie primitivis, $(a,b,c)$, $(a',b',c')$, $(a'',b'',c'')$ etc. eiusdem determinantis $D$, quarum termini antecedentes $a$, $a'$, $a''$ etc. sunt numeri inter se primi, forma $(A,B,C)$ ex illis omnibus composita invenitur. statuendo $A$ aequalem producto ex omnibus $a$, $a'$, $a''$ etc.; $B$ congruum ipsis $b$, $b'$, $b''$ etc. secundum modulos $a$, $a'$, $a''$ etc. resp.; $C=\frac{BB-D}{A}$. Facile enim perspicietur, ex duabus formis $(a,b,c)$, $(a',b',c')$ compositam fore formam $(aa',B,\frac{BB-D}{aa'})$; ex hac atque $(a'',b'',c'')$ formam $(aa'a'',B,\frac{BB-D}{aa'a''})$ etc. --- Vice versa

4) Proposita forma proprie primitiva $(A,B,C)$ determinantis $D$, si terminus $A$ in factores quotcunque inter se primos $a$, $a'$, $a''$ etc. resolvitur; numeri $b$, $b'$, $b''$ etc. ipsi $B$ vel aequales vel saltem sec. mod. $a$, $a'$, $a''$ etc. resp. congrui accipiuntur, atque fit $ac=bb-D$, $a'c'=b'b'-D$, $a''c''=b''b''-D$ etc.: forma $(A,B,C)$ composita erit e formis $(a,b,c)$, $(a'\ b'\ c')$, $(a''\ b''\ c'')$, sive \textit{in has formas resolubilis}. Nullo negotio probatur, eandem propositionem adhucdum valere, etiamsi forma $(A,B,C)$ sit improprie primitiva vel derivata. Hoc itaque modo quaelibet forma in alias eiusdem determinantis resolvi potest, quarum termini antecedentes omnes sint vel numeri primi vel numerorum primorum potestates. Talis resolutio saepenumero commode applicari potest, si ex pluribus formis datis componenda est una. Ita \textit{e. g.}\ si quaeritur forma composita e formis $(3,1,134)$, $(10,3,41)$, $(15,2,27)$, resolvatur secunda in has $(2,1,201)$, $(5,-2,81)$, tertia in has $(3,-1,134)$, $(5,2,81)$, patetque, formam ex quinque formis $(3,1,134)$, $(2,1,201)$, $(5,-2,81)$, $(3,-1,134)$, $(5,2,81)$ compositam, quocunque ordine accipiantur, etiam ex tribus datis compositam fore. At ex compositione primae cum quarta oritur forma principalis $(1,0,401)$; eadem provenit ex compositione tertiae cum quinta; quare ex compositione cunctarum conflatur forma $(2,1,201)$.

\SourceLeaf{270}
5) Propter rei utilitatem operae pretium est, hanc methodum adhuc amplius explicare. Ex observatione praecedente manifestum est, problema, quotcunque formas datas proprie primitivas eiusdem determinantis componere, reduci posse ad compositionem formarum, quarum termini initiales sint potestates numerorum primorum (nam numerus primus tamquam sui ipsius potestas prima considerari potest). Quamobrem eum imprimis casum contemplari convenit, ubi duae formae proprie primitivae $(a,b,c)$, $(a',b',c')$ sunt componendae, in quibus $a$ et $a'$ sunt potestates \textit{eiusdem} numeri primi. Sit itaque $a=h^{\varkappa}$, $a'=h^{\lambda}$ designante $h$ numerum primum, supponamusque (quod licet), $\varkappa$ non esse minorem quam $\lambda$. Erit itaque $h^{\lambda}$ div. comm. max. numerorum $a$, $a'$, qui si insuper ipsum $b+b'$ metitur, habebitur casus initio huius art. consideratus, eritque $(A,B,C)$ ex propositis composita si statuitur $A=h^{\varkappa-\lambda}$, $B\equiv b\ (\mathrm{mod.}\ h^{\varkappa-\lambda})$ et $\equiv b'\ (\mathrm{mod.}\ 1)$, quae conditio posterior manifesto omitti potest; $C=\frac{BB-D}{A}$. Si vero $h^{\lambda}$ ipsum $b+b'$ non metitur, necessario div. comm. max. horum numerorum et ipse erit potestas ipsius $h$, sit igitur $=h^{\nu}$, eritque $\nu<\lambda$ (statui debet $\nu=0$. si forte $h^{\lambda}$ et $b+b'$ inter se primi sunt). Si itaque $\mathfrak{P}'$, $\mathfrak{P}''$, $\mathfrak{P}'''$ ita determinantur, ut fiat
\[
\mathfrak{P}'h^{\varkappa}+\mathfrak{P}''h^{\lambda}+\mathfrak{P}'''(b+b')=h^{\nu}
\]
$\mathfrak{P}$ vero ad libitum assumitur, forma $(A,B,C)$ ex datis erit composita, si statuitur
\[
\begin{gathered}A=h^{\varkappa+\lambda-2\nu},\\
B=b+h^{\varkappa-\nu}(\mathfrak{P}h^{\lambda}-\mathfrak{P}'(b-b')-\mathfrak{P}'''c),\\
C=\frac{BB-D}{A}\end{gathered}
\]
Sed facile perspicitur, in hoc casu etiam $\mathfrak{P}'$ ad libitum assumi posse, quare statuendo $\mathfrak{P}=\mathfrak{P}'=0$, fit
\[
B=b-\mathfrak{P}'''ch^{\varkappa-\nu}
\]
sive generaliter
\[
B=kA+b-\mathfrak{P}'''ch^{\varkappa-\nu}
\]
designante $k$ numerum arbitrarium (art. praec.). In hanc formulam simplicissimam solus $\mathfrak{P}'''$ ingreditur, qui est valor expr. $\frac{h^{\nu}}{b+b'}\ (\mathrm{mod.}\ h^{\lambda})$\HistoricalNote{*)}{sive expr. $\frac{1}{\frac{b+b'}{h^{\nu}}}\ (\mathrm{mod.}\ h^{\lambda-\nu})$ unde $B\equiv b-\frac{ch^{\varkappa-\nu}}{\frac{b+b'}{h^{\nu}}}\equiv\frac{(D+bb'):h^{\nu}}{(b+b'):h^{\nu}}\ (\mathrm{mod.}\ A)$}. Si \textit{e. g.}\ quaeritur forma composita ex $(16,3,19)$ et $(8,1,37)$, est $h=2$, $\varkappa=4$, $\lambda=3$, $\nu=2$.
\SourceLeaf{271}
Hinc $A=8$, $\mathfrak{P}'''$ valor expr. $\frac{4}{4}\ (\mathrm{mod.}\ 8)$, qualis est $1$, unde $B=8k-73$, adeoque faciendo $k=9$, $B=-1$ atque $C=37$, sive $(8,-1,37)$ forma quaesita.

Propositis itaque formis quotcunque, quarum termini initiales omnes sunt potestates numerorum primorum, circumspiciendum erit, num aliquarum termini antecedentes sint potestates \textit{eiusdem} numeri primi, atque hae inter se respective per regulam modo traditam componendae. Hac ratione prodibunt formae, quarum termini primi etiamnum erunt potestates numerorum primorum, sed omnino diversorum; forma itaque ex his composita per observ. tertiam definiri poterit. \textit{E. g.}\ propositis formis $(3,1,47)$, $(4,0,35)$, $(5,0,28)$, $(16,2,9)$, $(9,7,21)$, $(16,6,11)$, ex prima et quinta conflatur forma $(27,7,7)$; ex secunda et quarta confit $(16,-6,11)$, ex hac et sexta $(1,0,140)$, quae negligi potest. Supersunt itaque $(5,0,28)$, $(27,7,7)$, ex quibus producitur $(135,-20,4)$, cuius loco assumi potest proprie aequivalens $(4,0,35)$. Haec itaque est resultans ex compositione sex propositarum.

Ceterum ex hoc fonte plura alia artificia in applicatione utilia hauriri possunt; sed ne nimis longi fiamus, uberiorem huius rei tractationem supprimimus, ad alia difficiliora properantes.

\Needspace{8\baselineskip}
\Article{244}
Si per formam aliquam $f$ repraesentari potest numerus $a$, per formam $f'$ numerus $a'$, atque forma $F$ in $ff'$ est transformabilis: nullo negotio perspicitur, productum $aa'$ per formam $F$ repraesentabile fore. Hinc statim sequitur, quando determinantes harum formarum sint negativi, formam $F$ positivam fore, si vel utraque $f$, $f'$ sit positiva vel utraque negativa; contra $F$ fieri negativam, si altera formarum $f$, $f'$ sit positiva altera negativa. Subsistamus in eo imprimis casu, quem in art. praec. consideravimus, ubi $F$ ex $f$, $f'$ composita est. atque $f$, $f'$ et $F$ eundem determinantem $D$ habent. Supponamus insuper, repraesentationes numerorum $a$, $a'$ per formas $f$, $f'$ fieri per valores indeterminatarum inter se primos, atque priorem pertinere ad valorem $b$ expressionis $\surd D\ (\mathrm{mod.}\ a)$, posteriorem ad valorem $b'$ expr. $\surd D\ (\mathrm{mod.}\ a')$, ponaturque $bb-D=ac$, $b'b'-D=a'c'$. Tunc per art. 168 formae $(a,b,c)$, $(a',b',c')$ proprie aequivalebunt formis $f$, $f'$; quare $F$ etiam ex illis duabus formis composita erit. Sed ex iisdem formis composita erit forma $(A,B,C)$, si, posito numerorum $a$, $a'$, $b+b'$ divisore communi maximo $=\mu$, statuitur $A=\frac{aa'}{\mu\mu}$, $B\equiv b$ et $b'$ sec. modulos $\frac{a}{\mu}$, $\frac{a'}{\mu}$. resp.,
